% Finite-dimensional nuclear endpoint results compiled on 26 October 2026.
% Evidence: SchattenNuclearEndpoint, RealNuclearMatching,
% RealNuclearFieldEquality, NuclearFieldComparison, RealNuclearOptimality,
% and SchattenNuclearCorrection. No all-optimizer or endpoint local-minimum
% classification is inferred from the strict-exponent Schatten theorems.
\section{The nuclear endpoint and real-field matching}
\label{sec:nuclear-endpoint}

For equal-rank subspaces $U,V$ of a finite-dimensional real or complex
Hilbert space, write
\[
 d_{\mathbb F}(U,V)
   =\min_{\substack{W\in\mathcal U_{\mathbb F}(E)\\ W(U)=V}}
       \|I-W\|_1,
 \qquad \mathbb F\in\{\mathbb R,\mathbb C\},
\]
where $\mathcal U_{\mathbb R}(E)$ is the orthogonal group and
$\mathcal U_{\mathbb C}(E)$ is the unitary group. The minimum exists in
both fields. We use the principal angles
$\theta_1,\ldots,\theta_k\in[0,\pi/2]$ with their original labels;
repeated values count separately.

Over $\mathbb C$, the optimum is
\begin{equation}
 d_{\mathbb C}(U,V)=2\sum_{i=1}^k\sin\theta_i
                  =\|P_U-P_V\|_1.
 \label{eq:nuclear-complex-endpoint}
\end{equation}
The all-pairs identity and attainment were established in
\cref{thm:complex-trace-all-pairs}. A second checked proof obtains the
nuclear lower bound by taking $p\downarrow1$ in the unrestricted complex
Schatten inequality of \cref{thm:schatten-independent-planes}. Continuity
of the finite singular-value power sum holds even at zero singular values;
a principal-plane phase construction attains the limiting value.

The real optimum has an exact finite combinatorial description.
Let $J=\{i:\theta_i>0\}$, retaining each angle label. Define
\begin{equation}
 s(\theta)=\min\{4\sin(\theta/2),\,2\},\qquad
 e(\theta,\psi)=4\sin\!\left(\frac{\theta+\psi}{2}\right).
 \label{eq:nuclear-matching-weights}
\end{equation}
An involution $\sigma:J\to J$ consists of fixed vertices and
pairwise exchanges. Its cost is
\begin{equation}
 C_1(\sigma)=
 \sum_{\substack{i\in J\\\sigma(i)=i}}s(\theta_i)
 +\frac12\sum_{\substack{i\in J\\\sigma(i)\ne i}}
          e(\theta_i,\theta_{\sigma(i)}).
 \label{eq:nuclear-matching-cost}
\end{equation}
The factor one half counts each exchanged pair once.

\begin{theorem}[Unrestricted real nuclear transport]
\label{thm:nuclear-real-matching}
For every pair of equal-rank subspaces in a finite-dimensional real
Hilbert space,
\begin{equation}
 \boxed{d_{\mathbb R}(U,V)=\min_{\sigma^2=\mathrm{id}}C_1(\sigma).}
 \label{eq:nuclear-real-minimum}
\end{equation}
An ambient orthogonal transporter attains the minimum. The formula
includes both determinant classes, all zero, repeated and right
principal angles, and arbitrary common ambient directions.
\end{theorem}

\begin{lstlisting}[style=leanpaper,
 caption={The checked real nuclear attainment theorem. The omitted context
 supplies the finite-dimensional real inner-product structure. The
 \texttt{realNuclearMatchingMinimum} takes the original strictly positive
 labelled angles.},
 label={lst:nuclear-real-matching}]
theorem exists_minimizing_real_nuclear_transport_matching
    (K L : Submodule Real E)
    (hKL : Module.finrank Real K = Module.finrank Real L) :
    ∃ T : LinearIsometryEquiv (RingHom.id Real) E E,
      K.map T.toLinearMap = L ∧
      UnitarilyInvariantSeminorm.schatten Real E E 1
        (LinearMap.id - T.toLinearMap) =
        realNuclearMatchingMinimum
          (fun j : {j : Fin (Module.finrank Real K) //
            0 < K.principalAngles L j} => K.principalAngles L j.val) ∧
      ∀ V : LinearIsometryEquiv (RingHom.id Real) E E,
        K.map V.toLinearMap = L ->
          UnitarilyInvariantSeminorm.schatten Real E E 1
            (LinearMap.id - T.toLinearMap) <=
          UnitarilyInvariantSeminorm.schatten Real E E 1
            (LinearMap.id - V.toLinearMap)
\end{lstlisting}

\noindent The displayed Lean signature suppresses typeclass context and
uses \texttt{Real} for the scalar field and the expanded isometry type;
the source theorem states the same attainment and comparison.
The proof passes the $1<p<2$ unrestricted real matching theorem to
$p=1$: both the minimum over the compact ambient transport fiber and
the finite matching minimum are continuous at the endpoint.
The singleton weight compares rotation with reflection in a real
principal plane. A matched pair couples two such planes in four real
dimensions. Thus a matching optimizes over \emph{all} ambient orthogonal
transports, rather than only block-preserving candidates.

\begin{theorem}[Exact real--complex equality criterion]
\label{thm:nuclear-real-complex-parity}
Take a real and a complex equal-rank subspace pair with the same
principal-angle list, counted with multiplicities. Their nuclear
transport minima satisfy $d_{\mathbb C}\le d_{\mathbb R}$, with equality
if and only if every angle $t\in(0,\pi/2)$ has even multiplicity:
\begin{equation}
 d_{\mathbb R}=d_{\mathbb C}
 \quad\Longleftrightarrow\quad
 \#\{i:\theta_i=t\}\ \text{is even for every }0<t<\pi/2.
 \label{eq:nuclear-parity}
\end{equation}
Zero angles and right angles impose no parity restriction.
The real and complex ambient dimensions need not agree; equality of
the accepted principal-angle data is an explicit hypothesis.
\end{theorem}

\begin{proof}
For $\theta\in[0,\pi/2]$, the singleton cost obeys
$s(\theta)\ge2\sin\theta$, with equality only at zero or a right
angle. Likewise,
$e(\theta,\psi)\ge2\sin\theta+2\sin\psi$,
with equality precisely for equal angles. Every matching therefore
costs at least $2\sum_i\sin\theta_i$.
Equality requires all strictly acute positive angles to be paired
with an identical angle. Such a fixed-point-free pairing exists
within every strictly acute angle level exactly when that level has
even cardinality. The remaining zero and right-angle labels can be
left unmatched. The unrestricted matching theorem and
\cref{eq:nuclear-complex-endpoint} give the claim.
\end{proof}

\begin{theorem}[When the real minimum equals the direct-rotation value]
\label{thm:nuclear-real-canonical-value}
For an equal-rank real pair with principal-angle list
$\theta_1,\ldots,\theta_k$,
\begin{equation}
 d_{\mathbb R}(U,V)=4\sum_{i=1}^k\sin(\theta_i/2)
 \quad\Longleftrightarrow\quad
 \left\{\begin{array}{l}
   \theta_i\le\pi/3\ \text{for every }i,\\
   \#\{i:\theta_i>0\}\le1.
 \end{array}\right.
 \label{eq:nuclear-real-canonical-value}
\end{equation}
\end{theorem}

\begin{proof}
If the matching minimum equals the canonical sum, the all-singleton
matching forces $s(\theta_i)=4\sin(\theta_i/2)$ for every label,
which is equivalent to $\theta_i\le\pi/3$. For any two positive
angles, pairing their vertices gives a strict saving over their two
canonical rotation costs, excluding a second positive label.
Conversely, with at most one positive angle at most $\pi/3$,
no exchange of distinct positive labels is possible, and the
canonical value is attained.
\end{proof}

The last theorem is stated in Lean as an equivalence for the
\emph{minimum value}. It does not by itself classify the full set of
ambient real optimizers. It also shows why a one-plane reflection cutoff
cannot establish the original Proposition~4.4: two equal
$30^\circ$ angles give
\[
  d_{\mathbb R}=4\sin(30^\circ)=2
  \quad<\quad8\sin(15^\circ)\approx2.07055,
\]
where the right side is the direct-rotation nuclear displacement.
For a single $30^\circ$ angle, by contrast, the real minimum is
$4\sin(15^\circ)>2\sin(30^\circ)=d_{\mathbb C}$.

\begin{corollary}[Nearest subspace-preserving complex unitary]
\label{cor:nuclear-unitary-rounding}
Let $A$ be unitary on a finite-dimensional complex Hilbert space, and
let $K$ be any subspace. Among all unitaries $V$ preserving $K$,
there is an attainer with
\begin{equation}
 \boxed{\displaystyle
 \min_{\substack{V\ \mathrm{unitary}\\V(K)=K}}\|A-V\|_1
 =2\sum_{j=1}^{\dim K}\sin\theta_j(AK,K).}
 \label{eq:nuclear-unitary-rounding}
\end{equation}
\end{corollary}

\begin{proof}
The bijection between transporters of $AK$ onto $K$ and corrections
preserving $K$ is given by $T=VA^{-1}$. Right-unitary invariance gives
$\|A-V\|_1=\|I-VA^{-1}\|_1$. Apply the complex nuclear transport
minimum to the pair $AK,K$.
\end{proof}

The real analogue is also attained, with a different value formula.
For an orthogonal $A$ and a real subspace $K$, let
$L=A(K)$ and let $\theta_1,\ldots,\theta_k$ be the principal angles
of $L,K$. The minimum of $\|A-V\|_1$ over all orthogonal $V$ with
$V(K)=K$ equals the matching minimum
$\min_{\sigma^2=\mathrm{id}} C_1(\sigma)$ of
\cref{eq:nuclear-matching-cost} for this angle list. This value is
attained by an orthogonal $V$ commuting with $P_K$, and the comparison
includes every orthogonal commuting competitor. The corresponding theorems are in the checked Lean modules
\texttt{RealSchattenCommutingCorrection} and
\texttt{SchattenNuclearCorrection}.

For $1<p<2$, the complex all-optimizer and local-to-global theorems in
\cref{sec:schatten-optimization} remain available. At $p=1$, the
results above establish exact values and attainment. They do not assert
that every complex nuclear optimizer has been classified or that every
constrained nuclear local minimum is global. The parity theorem is an
exact endpoint result; its proposed analogue for the real/complex
strict-exponent range is a separate question.
